%=====================================================================
\section{Axial march and the melt front}
\label{sec:front}

%---------------------------------------------------------------------
\subsection{Effectiveness--NTU segment march}
\label{sec:front:march}

The tube is divided into $n_{s}$ segments of length $\Delta z =
L_{\mathrm{tube}}/n_{s}$. Nodes are placed at $z_{j}=j\Delta z$, $j=0,\dots,n_s$,
and each node is assigned the melting temperature of the PCM layer containing
it,
\begin{equation}
  \Tm^{(j)} \;=\; T_{m,\ell}
  \quad\text{for}\quad
  z_{j} \in \bigl[\,z_{\ell},\, z_{\ell+1}\bigr),
  \qquad
  z_{\ell} = \frac{\ell}{N_{\mathrm{lay}}}\,L_{\mathrm{tube}} .
\end{equation}
During charging the layer melting temperatures descend with the fluid glide,
$T_{m,\ell}=\Tm - \ell\,\Delta T_{3C,2C}/N_{\mathrm{lay}}$; during discharging
they ascend, $T_{m,\ell}^{\mathrm{dc}}=\Tm - \Delta T_{3C,2C} +
(\ell+1)\Delta T_{3C,2C}/N_{\mathrm{lay}}$.

Within a segment the PCM surface is isothermal at $\Tm^{(j)}$, so the segment is
a single-stream heat exchanger of effectiveness $1-e^{-\mathrm{NTU}}$:
\begin{equation}
  \mathrm{NTU}_{j}
  \;=\;
  \frac{2\pi r_{i}\,\Ui^{(j)}\,\Delta z}{\md\,c_{p}},
  \qquad
  T_{j+1}
  \;=\;
  \Tm^{(j)} + \bigl(T_{j}-\Tm^{(j)}\bigr)\,e^{-\mathrm{NTU}_{j}} .
  \label{eq:ntu}
\end{equation}
The heat given up by the fluid over the segment, per unit tube length, is
\begin{equation}
  q'_{j} \;=\; \frac{\md\,c_{p}\,\bigl(T_{j}-T_{j+1}\bigr)}{\Delta z},
  \label{eq:qprime}
\end{equation}
positive when the fluid is melting PCM and negative when the PCM is freezing.
The march is sequential: $T_{0}=T_{\mathrm{inlet}}$, and each segment's outlet is
the next segment's inlet.

\emph{Source:} \code{\_legacy\_physics.temperature\_profile\_melt};
\code{well\_marched.march}.

\medskip
\noindent\textbf{The two formulations report different heat rates.}
Formulation~A forms the tube heat rate as an \emph{endpoint enthalpy
difference} over the whole tube,
\begin{equation}
  Q_{\mathrm{A}} \;=\; \md\,\bigl[h(T_{0},P) - h(T_{n_{s}},P)\bigr],
  \label{eq:Q-endpoint}
\end{equation}
evaluated by CoolProp, whereas Formulation~B accumulates the segment
contributions actually taken up by the PCM,
$Q_{\mathrm{B}} = \sum_{j} q'_{\mathrm{eff},j}\,\Delta z$
(Eq.~\eqref{eq:qeff}).

Equation~\eqref{eq:Q-endpoint} is the better measure of what the \emph{fluid}
gave up, since it uses real enthalpies rather than the local $c_{p}$ of
Eq.~\eqref{eq:qprime}. But it is not the heat the \emph{front} absorbed, and in
Formulation~A nothing requires the two to agree --- the front is advanced from
$\Tre$ through Eq.~\eqref{eq:stefan-explicit}, never from $Q_{\mathrm{A}}$.
That is the closure defect of \S\ref{sec:defects} stated in one line.

%---------------------------------------------------------------------
\subsection{Formulation A --- closed-form Stefan front (v0.1)}
\label{sec:front:legacy}

This is the formulation behind the conference paper. It is selected by
\code{Case(front="closed\_form")} and is retained as a regression reference.

The front is obtained from a quasi-steady Stefan solution for a bare concentric
cylinder, driven by the outer-surface temperature $\Tre$ of
Eq.~\eqref{eq:Tre}. With $\alpha_{m}=\km/(\rhom c_{p,m})$, the Fourier and
phase-change numbers
\begin{equation}
  \mathrm{Fo} \;=\; \frac{\alpha_{m}t}{r_{e}^{2}},
  \qquad
  \mathrm{Ph} \;=\;
  \left|\frac{\hm}{c_{p,m}\bigl(\Tre-\Tm\bigr)}\right|
  \;=\; \frac{1}{\mathrm{Ste}},
\end{equation}
and the dimensionless front position $x = 1 + \delta/r_{e}$, the front solves
\begin{equation}
  \boxed{\;
  \mathrm{Fo} \;=\; \mathrm{Ph}\;\Bigl[
    \tfrac{1}{2}x^{2}\ln x - \tfrac{1}{4}x^{2} + \tfrac{1}{4}
  \Bigr]
  \;}
  \label{eq:stefan-implicit}
\end{equation}
which, on substituting $\mathrm{Fo}$ and $\mathrm{Ph}$, is equivalently
\begin{equation}
  \tfrac{1}{2}x^{2}\ln x - \tfrac{1}{4}x^{2} + \tfrac{1}{4}
  \;=\;
  \frac{\km\bigl(\Tre-\Tm\bigr)\,t}{\rhom\,\hm\,r_{e}^{2}} .
  \label{eq:stefan-explicit}
\end{equation}
Equation~\eqref{eq:stefan-explicit} is the exact quadrature of
$\rhom\hm\,\mathrm{d}(\pi s^{2})/\mathrm{d}t = 2\pi\km(\Tre-\Tm)/\ln(s/r_{e})$,
with $s=r_{e}+\delta$, and is solved for $\delta$ by safeguarded
Newton--bisection on $[0,\delta_{\max}]$.

\emph{Source:} \code{\_legacy\_physics.compute\_delta2\_fast}.

\medskip
\noindent\textbf{Two assumptions, both violated.} Equation~\eqref{eq:stefan-explicit}
is derived for a \emph{bare} cylinder of perimeter $2\pi r_{e}$, whereas the
fluid side gives up its heat through the finned perimeter $P_{T}$ of
Eq.~\eqref{eq:perim}. The two surfaces differ by $3.72$, so the heat leaving the
fluid and the latent heat absorbed by the front are not the same energy. Second,
it assumes $\Tre$ has held constant since $t=0$, whereas $\Tre$ is applied at its
current value; measured at mid-well, $\Tre-\Tm$ runs from $0$ to
\SI{8.62}{\kelvin} across the charging window.

These two errors act in opposite directions and are quantified in
\S\ref{sec:defects}. Neither is a numerical defect: the solver returns the exact
root of Eq.~\eqref{eq:stefan-explicit} to a residual of $10^{-15}$. It is the
premises that fail.

%---------------------------------------------------------------------
\subsection{Formulation B --- energy-balance front (v0.2)}
\label{sec:front:marched}

The correction observes that of the two models --- heat delivered, and front
position --- only one may be specified freely. The other follows from
conservation. The front is therefore advanced by the heat the network actually
delivers:
\begin{equation}
  \boxed{\;
  \rhom\,\hm\,\frac{\mathrm{d}\Amelt}{\mathrm{d}t} \;=\; q'(t) .
  \;}
  \label{eq:marched}
\end{equation}

Integrated explicitly over the time levels $t^{n}$, with $\Delta t = t^{n+1}-t^{n}$,
\begin{equation}
  \Amelt^{n+1}
  \;=\;
  \mathcal{P}\!\left[\,\Amelt^{n} + \frac{q'\,\Delta t}{\rhom \hm}\,\right],
  \qquad
  \mathcal{P}[A] = \min\bigl(\max(A,0),\,A_{\max}\bigr),
  \label{eq:advance}
\end{equation}
where the projection $\mathcal{P}$ enforces two physical bounds: PCM that is
already solid cannot freeze further, and a segment cannot melt more PCM than it
contains. The heat actually taken up is recovered from the projected state,
\begin{equation}
  q'_{\mathrm{eff}}
  \;=\;
  \frac{\bigl(\Amelt^{n+1}-\Amelt^{n}\bigr)\rhom \hm}{\Delta t},
  \label{eq:qeff}
\end{equation}
and it is $q'_{\mathrm{eff}}$, not $q'$, that is accumulated into the delivered
energy and that sets the segment outlet temperature in Eq.~\eqref{eq:ntu}.
Accumulating $q'$ instead would credit the fluid with heat no PCM supplied.

Three properties follow, none of which Formulation A possesses.

\begin{description}[leftmargin=1.6em,itemsep=2pt]
\item[Closure is structural.] $\int q'\,\mathrm{d}t = \rhom\hm\Delta\Vmelt$
  holds identically rather than approximately. This is a consistency property,
  \emph{not} evidence that the model is correct: Eqs.~\eqref{eq:advance}
  and~\eqref{eq:qeff} are inverse operations, so a residual of $10^{-16}$
  measures floating-point arithmetic and nothing else.
\item[The fins enter the front automatically,] because $q'$ comes from the
  network of Eq.~\eqref{eq:Ui}, which carries $P_{T}$ and $\etao$.
\item[$\Tre$ need not be constant.] The state is marched, so a varying inlet
  temperature is handled correctly and Eq.~\eqref{eq:Tre} is not required.
\end{description}

%---------------------------------------------------------------------
\subsection{The fins are metal, not PCM}
\label{sec:front:finmetal}

$\Amelt$ is the quantity the latent balance conserves, so it must be PCM and
nothing else. The annulus between $r_{e}$ and $r_{e}+\delta$ is not all PCM: the
fins occupy metal inside it, of cross-section $n_{f}t_{f}\min(\delta,L_{f})$,
which saturates once the front passes the fin tips. Hence
\begin{equation}
  \Amelt \;=\; \pi\left[(r_{e}+\delta)^{2}-r_{e}^{2}\right]
  \;-\; n_{f}\,t_{f}\,\min\!\left(\delta,\,L_{f}\right).
  \label{eq:area-fins}
\end{equation}

Inverting Eq.~\eqref{eq:area-fins} for $\delta$ is still closed form --- the fin
term only shifts the quadratic, and each branch is solved exactly:
\begin{align}
  \delta \le L_{f}: &\qquad
    \pi\delta^{2}+\left(2\pi r_{e}-n_{f}t_{f}\right)\delta-\Amelt = 0,
  \label{eq:delta-in}\\
  \delta > L_{f}: &\qquad
    \pi\delta^{2}+2\pi r_{e}\delta-\left(\Amelt+n_{f}t_{f}L_{f}\right) = 0 .
  \label{eq:delta-out}
\end{align}
With $n_{f}=0$ both collapse to $\delta=\sqrt{r_{e}^{2}+\Amelt/\pi}-r_{e}$, the
bare concentric annulus used before this correction.

\medskip
\noindent\textbf{Corrected in v0.3.} Earlier versions treated the whole annulus
as PCM. That overstated the melted volume by \SI{11}{\percent} when the layer is
thin and \SI{6}{\percent} at the design point, and understated $\delta$ ---
hence the melt-layer conduction resistance of Eq.~\eqref{eq:he} --- by up to
\SI{34}{\percent}. Because it inflated $\varepsilon_{\mathrm{PCM}}$ it also
pushed $N_{\mathrm{inventory}}$ slightly \emph{up}; correcting it moves the
design point from $12.79$ to $12.74$ wells.

What remains unmodelled is the front \emph{shape}. The fins enhance the heat
transfer through $\etao$, but the melt region is still taken as a circular
annulus (hypothesis H3); physically it bulges along the fins.

%---------------------------------------------------------------------
\subsection{Formulation C --- enthalpy state (v0.4)}
\label{sec:front:enthalpy}

\emph{This subsection is written at greater length than the rest of the
document. Formulation~C is the part of the model a new student is most likely to
have to modify, and it is the part where the reasoning --- not the final
formula --- is what has to be understood. Readers who only want the equations
can take Eqs.~\eqref{eq:enth-curve}, \eqref{eq:enth-ode} and
\eqref{eq:enth-exact} and skip the discussion.}

\subsubsection{The question a state variable answers}

Every transient model must decide what it \emph{remembers} from one instant to
the next. That choice --- the state variable --- is not cosmetic. It decides
which physical situations the model can represent at all, because anything the
state cannot encode is a situation the model cannot be in.

The three formulations in this section are three answers to that one question,
and they are best read as a sequence:

\begin{center}
\begin{tabular}{@{}l l l@{}}
\toprule
 & state variable & what it can represent \\
\midrule
A (v0.1) & $\delta$, from a closed form & a melt front, at one prescribed $\Tre$ \\
B (v0.2) & $\Amelt$, integrated & a melt front, driven by the actual heat flow \\
C (v0.4) & $E'$, integrated & a melt front \emph{and} a PCM temperature \\
\bottomrule
\end{tabular}
\end{center}

\noindent
Formulations A and B both store \emph{how much has melted} and nothing else.
Ask either of them ``how hot is the PCM in segment~$j$?'' and the only available
answer is $\Tm$ --- because that is all the state can say. While some solid
remains that answer is correct, and it is why the omission went unnoticed for so
long.

\subsubsection{Where Formulation B breaks, and why it is not a small error}

Consider a segment near the well head during charging. The secondary fluid
enters at $T_{4c}=\SI{160}{\celsius}$; the PCM melts at
$\Tm=\SI{150}{\celsius}$. That segment sees the hottest fluid, so it melts
first. Once its last solid is gone, the resistance network still reports a
positive heat flow --- the fluid is still \SI{10}{\kelvin} hotter than the
melting point, and nothing in Eq.~\eqref{eq:Ui} knows the segment is finished.

Formulation~B must then do one of two things, and both are wrong:

\begin{enumerate}[leftmargin=1.7em,itemsep=2pt]
\item \textbf{Keep melting.} $\Amelt$ grows past the PCM the segment actually
  contains. This is what the unguarded version did, and it produced local melt
  fractions up to $\varepsilon_{\mathrm{local}}=1.78$: the model drew latent
  heat from material that was not there.
\item \textbf{Clip and discard.} Cap $\Amelt$ at $A_{\mathrm{avail}}$ and throw
  the surplus away. Now the melt fraction is physical, but energy is not
  conserved --- and the fluid leaves the segment at the wrong temperature,
  because it gave up heat that went nowhere.
\end{enumerate}

The second is what the guarded version did, and the magnitude is the reason this
is not a rounding-level concern. At the design point the heat discarded during
\emph{discharge} reached \SI{148}{\percent} of the heat delivered. Most of what
the resistance network asked for was thrown away, and the round-trip figures
were computed on what was left.

Notice that no amount of care inside the march can fix this. The defect is in
the \emph{state}: a variable that counts melted area has no slot for ``liquid
above $\Tm$''. The physics has to be given somewhere to live.

\subsubsection{Why enthalpy, and not temperature}

The obvious repair is to carry a PCM temperature alongside the melted area. It
is also the wrong repair, and the reason is worth understanding because it
recurs throughout phase-change modelling.

At a phase change, temperature is a \emph{bad} state variable: it is not
invertible. During melting the PCM sits at $\Tm$ while its energy content
climbs by the entire latent heat, so $T$ carries no information at all
throughout the process that dominates the problem. Carrying $(T, \Amelt)$ as a
pair works, but then the code must decide at every step which of the two is
active, keep them consistent, and handle the moments when the active variable
changes.

Enthalpy is invertible. Every state of the PCM --- subcooled solid, mid-melt,
superheated liquid --- corresponds to exactly one value of $E'$, and $E'$ moves
monotonically as heat is added. So a \emph{single} scalar per segment encodes
all three regimes, and the regime is read off from the value rather than tracked
separately. Figure~\ref{fig:enthalpy-curve} is the whole idea in one picture.

\begin{figure}[htbp]
\centering
\begin{tikzpicture}[>=Stealth, scale=1.0]
  % axes
  \draw[->,thick] (0,0) -- (9.6,0) node[right] {$T_{\mathrm{pcm}}$};
  \draw[->,thick] (0,0) -- (0,5.8) node[above] {$E'$};
  % melting temperature
  \draw[dashed,gray] (4.9,0) -- (4.9,5.4);
  \node[below] at (4.9,0) {$\Tm$};
  % latent plateau markers
  \draw[dashed,gray] (0,1.5) -- (4.9,1.5) node[pos=0,left] {$0$};
  \draw[dashed,gray] (0,3.5) -- (4.9,3.5) node[pos=0,left] {$E'_{\mathrm{lat}}$};
  % the curve: solid branch, plateau, liquid branch
  \draw[very thick] (1.2,0.05) -- (4.9,1.5);
  \draw[very thick] (4.9,1.5) -- (4.9,3.5);
  \draw[very thick] (4.9,3.5) -- (8.9,5.05);
  % branch labels, all kept clear of the curve and of each other
  \node[align=left,font=\small,anchor=west] at (0.85,2.65)
    {solid, subcooled\\[1pt] slope $C_{s}$};
  \node[align=center,font=\small] at (6.9,5.3)
    {liquid, superheated\\[1pt] slope $C_{l}$};
  % latent plateau annotation, placed to the LEFT of the vertical segment
  \draw[<->,thick,black!60] (4.55,1.5) -- (4.55,3.5);
  \node[font=\small,align=right,anchor=east,black!70] at (4.4,1.95)
    {latent heat\\$\rhom\hm A_{\mathrm{avail}}$};
  % melt-fraction callout, to the RIGHT and well below the liquid branch
  \filldraw (4.9,2.6) circle (1.6pt);
  \draw[->,black!60] (5.9,1.75) -- (4.99,2.53);
  \node[font=\small,align=left,anchor=west,black!70] at (5.95,1.6)
    {position here is\\$\Amelt = E'/(\rhom\hm)$};
\end{tikzpicture}
\caption{The enthalpy curve, Eq.~\eqref{eq:enth-curve}. Reading it as
$E'(T)$ it is multivalued at $\Tm$ --- which is exactly why $T$ cannot be the
state. Reading it as $T(E')$, the direction the code uses, it is single-valued
everywhere. Position on the vertical segment is the melt fraction.}
\label{fig:enthalpy-curve}
\end{figure}

\subsubsection{The state, its capacities, and their sizes}

Let $E'$ be the enthalpy per unit tube length of the PCM belonging to one
segment, measured from a datum of \emph{fully solid at $\Tm$}. The datum is
arbitrary but this choice puts the two branch boundaries at $0$ and
$E'_{\mathrm{lat}}$, which keeps the tests in the code short.

With $A_{\mathrm{avail}}=V_{\mathrm{well}}/(n_{t}L_{\mathrm{tube}})$ the PCM
cross-section a segment actually owns,
\begin{equation}
  E'_{\mathrm{lat}} = \rhom \hm A_{\mathrm{avail}},
  \qquad
  C_{s} = \rhom c_{p,s} A_{\mathrm{avail}},
  \qquad
  C_{l} = \rhom c_{p,l} A_{\mathrm{avail}},
  \label{eq:enth-caps}
\end{equation}
are the latent capacity and the two sensible capacities, all per unit tube
length. Their numerical values at the design point are worth committing to
memory, because they explain the behaviour of the whole formulation:

\begin{center}
\begin{tabular}{@{}l r l@{}}
\toprule
quantity & value & \\
\midrule
$A_{\mathrm{avail}}$      & \num{4.5411e-3} & \si{\meter\squared} \\
$E'_{\mathrm{lat}}$       & \num{2.675}     & \si{\mega\joule\per\meter} \\
$C_{s}$                   & \num{9.010e3}   & \si{\joule\per\meter\per\kelvin} \\
$C_{l}$                   & \num{1.267e4}   & \si{\joule\per\meter\per\kelvin} \\
\midrule
$C_{l}/E'_{\mathrm{lat}}$ & \num{4.74e-3}   & \si{\per\kelvin} \\
\bottomrule
\end{tabular}
\end{center}

\noindent
That last ratio is the Stefan number in disguise. One kelvin of superheat is
worth \SI{0.47}{\percent} of the latent capacity, so the sensible branches are
\emph{energetically} minor --- with the \SI{10}{\kelvin} approach available
here, at most about \SI{4.7}{\percent}. A reader could reasonably conclude that
the whole formulation is a small correction. It is not, and
\S\ref{sec:front:levelling} explains why: the sensible branches matter through
the \emph{driving temperature difference}, not through the energy they store.

The enthalpy curve and the melted area follow directly:
\begin{equation}
  T_{\mathrm{pcm}} =
  \begin{dcases}
    \Tm + E'/C_{s}, & E' < 0 \quad\text{(solid, subcooled)},\\[2pt]
    \Tm, & 0 \le E' \le E'_{\mathrm{lat}} \quad\text{(two-phase)},\\[2pt]
    \Tm + \bigl(E'-E'_{\mathrm{lat}}\bigr)/C_{l},
      & E' > E'_{\mathrm{lat}} \quad\text{(liquid, superheated)},
  \end{dcases}
  \qquad
  \Amelt = \frac{\min\!\left(\max(E',0),\,E'_{\mathrm{lat}}\right)}{\rhom \hm}.
  \label{eq:enth-curve}
\end{equation}

The clipping in $\Amelt$ is not a guard bolted on to prevent bad behaviour ---
it is the definition. Melted area saturates because a segment cannot melt PCM it
does not own, and superheat is where the extra energy goes. This is the sense in
which $\varepsilon_{\mathrm{local}}\in[0,1]$ is \emph{structural} in
Formulation~C: there is no code path that can violate it, so no test is needed.

\subsubsection{The segment equation, derived}
\label{sec:front:segeq}

Equation~\eqref{eq:enth-ode} below is the heart of the march, and it is short
enough to look like a definition. It is not: it is the result of writing an
energy balance on two different control volumes and then eliminating the
quantity they share. This subsection does that in full, because the conductance
$K$ that comes out is \emph{not} the $\Ui P_{i}$ a reader would write down
unaided, and the difference reaches \SI{33}{\percent} at this design point.

\paragraph{The two control volumes.}
Fix a segment $j$, spanning $z\in[z_{j},z_{j+1}]$ with $\Delta z =
z_{j+1}-z_{j}$. Two control volumes share the tube wall over that interval:

\begin{description}[leftmargin=1.9em,itemsep=2pt]
\item[CV\,1 --- the PCM] belonging to segment $j$: cross-section
  $A_{\mathrm{avail}}$, length $\Delta z$, stationary, closed.
\item[CV\,2 --- the fluid] inside the tube over the same interval:
  cross-section $A_{\mathrm{flow}}=\pi r_{i}^{2}$, open, with mass flow $\md$
  entering at $T_{j}$ and leaving at $T_{j+1}$.
\end{description}

\noindent
Write $q'_{j}$ for the heat per unit tube length crossing from CV\,2 into
CV\,1. It appears in both balances with opposite sign, and eliminating it is
the whole exercise.

\medskip
\noindent\textbf{Step 1 --- the PCM balance.}
CV\,1 is closed and stationary, so the first law reduces to a statement about
its enthalpy alone. With $E'$ the enthalpy per unit tube length
(\S\ref{sec:front:enthalpy}), the enthalpy of the control volume is
$E'\Delta z$ and
\begin{equation}
  \frac{\mathrm{d}}{\mathrm{d}t}\bigl(E'_{j}\,\Delta z\bigr)
  \;=\;
  \underbrace{q'_{j}\,\Delta z}_{\text{through the tube wall}}
  \;\underbrace{-\;\bigl[\dot Q_{\mathrm{ax}}\bigr]_{z_{j}}^{z_{j+1}}}
              _{\text{axial conduction, H2}}
  \;+\;\underbrace{\dot Q_{\mathrm{form}}}_{\text{formation, H6}}
  \;+\;\underbrace{\dot Q_{\mathrm{azi}}}_{\text{neighbours, H9}} ,
\end{equation}
and the last three terms are dropped by hypothesis. Two of those deletions are
worth a number rather than a citation:

\begin{itemize}[leftmargin=1.5em,itemsep=2pt]
\item \textbf{Axial conduction (H2)} is negligible by an enormous margin. The
  axial diffusion time over one segment is $\Delta z^{2}/\alpha_{l} =
  \SI{5.4e9}{\second}$ against a \SI{3.6e4}{\second} charging window --- a ratio
  of $1.5\times10^{5}$. With $\Delta z = \SI{30.5}{\meter}$ this is not a close
  call, and it would remain safe at a hundred times the segment count.
\item \textbf{Azimuthal conduction (H9)} is \emph{not} negligible and is
  assumed away pending the calculation described in \S\ref{sec:defects:h9}. At
  the wellhead the neighbouring cell is \SI{48.9}{\kelvin} away and a slab
  estimate gives $\sim\SI{22}{\watt\per\meter}$ against a useful duty of
  $\sim\SI{83}{\watt\per\meter}$. This is the least defensible term in the
  model.
\end{itemize}

\noindent
Dividing by $\Delta z$ leaves the PCM balance in the form used by the code:
\begin{equation}
  \boxed{\;\frac{\mathrm{d}E'_{j}}{\mathrm{d}t} \;=\; q'_{j}\;}
  \label{eq:pcm-balance}
\end{equation}

\medskip
\noindent\textbf{Step 2 --- the fluid balance.}
CV\,2 is open. For incompressible single-phase flow (H7) with no viscous
dissipation, the energy equation per unit tube length is
\begin{equation}
  \underbrace{\rho A_{\mathrm{flow}} c_{p}\,\frac{\partial T}{\partial t}}
             _{\text{fluid thermal inertia}}
  \;+\;
  \md\,c_{p}\,\frac{\partial T}{\partial z}
  \;=\;
  -\,q'(z) .
  \label{eq:fluid-pde}
\end{equation}

The model drops the first term, and this deserves justification rather than
assertion. Two comparisons, pointing the same way:

\begin{center}
\begin{tabular}{@{}l r l@{}}
\toprule
fluid transit time over $L_{\mathrm{tube}}$ & \SI{2264}{\second} & ($u=\SI{1.35}{\meter\per\second}$) \\
charging window $t_{\mathrm{ch}}$           & \SI{36000}{\second} & \\
\midrule
ratio                                        & \num{0.063} & \\
$\rho A_{\mathrm{flow}}c_{p} \,/\, C_{l}$    & \num{0.257} & fluid vs.\ PCM sensible capacity \\
$\rho A_{\mathrm{flow}}c_{p}\Delta T \,/\, E'_{\mathrm{lat}}$
   & \num{0.0122} & fluid swing vs.\ PCM latent, $\Delta T=\SI{10}{\kelvin}$ \\
\bottomrule
\end{tabular}
\end{center}

\noindent
The fluid re-establishes its profile in one transit, $\SI{2264}{\second}$,
while the PCM state evolves over \SI{10}{\hour}: a $16\times$ separation, so
quasi-steady is reasonable but \emph{not} asymptotic. The honest reading is
that the first \SI{38}{\minute} of each half-cycle is a startup transient the
model does not resolve, worth about \SI{6}{\percent} of the window. What makes
this tolerable is the third row: the energy parked in the fluid is only
\SI{1.2}{\percent} of the latent inventory it is charging, so the transient
misplaces very little energy even though it lasts a noticeable fraction of the
run.

Setting $\partial T/\partial t = 0$ and using the fact that CV\,1 presents a
\emph{single} temperature $T_{\mathrm{pcm}}$ to the whole segment
(H5, the lumped state of \S\ref{sec:front:enthalpy}), the driving difference is
$T(z)-T_{\mathrm{pcm}}$ and the local flux closes on the network of
\S\ref{sec:network}:
\begin{equation}
  q'(z) \;=\; \Ui\,\bigl(2\pi r_{i}\bigr)\,\bigl[T(z)-T_{\mathrm{pcm}}\bigr].
  \label{eq:local-flux}
\end{equation}
Note that $2\pi r_{i}$ appears, not $P_{T}$: $\Ui$ was referred to the
\emph{inner} area in Eq.~\eqref{eq:Ui}, and the finned perimeter is already
inside it through $R'_{\mathrm{outer}}$. Using $P_{T}$ here would count the fins
twice --- which is a compact statement of the v0.1 defect.

\medskip
\noindent\textbf{Step 3 --- integrate across the segment.}
Substituting Eq.~\eqref{eq:local-flux} into the steady form of
Eq.~\eqref{eq:fluid-pde} gives a linear ODE in $z$ with $T_{\mathrm{pcm}}$ and
$\Ui$ held constant over the segment:
\begin{equation}
  \md c_{p}\frac{\mathrm{d}T}{\mathrm{d}z}
  = -2\pi r_{i}\Ui\bigl(T-T_{\mathrm{pcm}}\bigr)
  \quad\Longrightarrow\quad
  \frac{\mathrm{d}\bigl(T-T_{\mathrm{pcm}}\bigr)}{T-T_{\mathrm{pcm}}}
  = -\frac{2\pi r_{i}\Ui}{\md c_{p}}\,\mathrm{d}z .
\end{equation}
Integrating from $z_{j}$ to $z_{j+1}$,
\begin{equation}
  T_{j+1} \;=\; T_{\mathrm{pcm}}
    + \bigl(T_{j}-T_{\mathrm{pcm}}\bigr)e^{-\mathrm{NTU}_{j}},
  \qquad
  \mathrm{NTU}_{j} \;=\; \frac{2\pi r_{i}\,\Ui\,\Delta z}{\md\,c_{p}},
  \label{eq:ntu-derived}
\end{equation}
which is Eq.~\eqref{eq:ntu}. It is the single-stream heat-exchanger result: one
stream with finite capacity rate $\md c_{p}$, the other an isothermal reservoir
of infinite capacity rate, effectiveness $\varepsilon = 1-e^{-\mathrm{NTU}}$.
The reservoir idealisation is exactly what the lumped PCM state licenses.

\medskip
\noindent\textbf{Step 4 --- the closure.}
The heat the fluid actually gives up over the segment is the enthalpy it loses,
\begin{equation}
  q'_{j}\,\Delta z \;=\; \md c_{p}\bigl(T_{j}-T_{j+1}\bigr)
  \;\overset{\eqref{eq:ntu-derived}}{=}\;
  \md c_{p}\bigl(T_{j}-T_{\mathrm{pcm}}\bigr)\bigl(1-e^{-\mathrm{NTU}_{j}}\bigr),
\end{equation}
and dividing by $\Delta z$ gives the conductance that closes
Eq.~\eqref{eq:pcm-balance}:
\begin{equation}
  \boxed{\;
  q'_{j} = K_{j}\bigl(T_{j}-T_{\mathrm{pcm}}\bigr),
  \qquad
  K_{j} = \frac{\md\,c_{p}\bigl(1-e^{-\mathrm{NTU}_{j}}\bigr)}{\Delta z} .
  \;}
  \label{eq:K-derived}
\end{equation}

Combining Eqs.~\eqref{eq:pcm-balance} and \eqref{eq:K-derived} with the
enthalpy curve $T_{\mathrm{pcm}}(E')$ of Eq.~\eqref{eq:enth-curve} closes the
system. This is the segment equation:
\begin{equation}
  \boxed{\;
  \frac{\mathrm{d}E'}{\mathrm{d}t}
  \;=\; K\left(T_{j} - T_{\mathrm{pcm}}(E')\right),
  \qquad
  K = \frac{\md\,c_{p}\left(1-e^{-\mathrm{NTU}}\right)}{\Delta z}.
  \;}
  \label{eq:enth-ode}
\end{equation}

One unknown per segment, $E'$; everything else is either geometry, a property,
or the upstream fluid temperature $T_{j}$ handed down by the segment before it.

\paragraph{Why $K$ is not $\Ui P_{i}$.}
This is the step most easily got wrong, so it is worth isolating. The naive
closure takes Eq.~\eqref{eq:local-flux} with $T(z)\approx T_{j}$, giving
$q'_{j}\approx 2\pi r_{i}\Ui (T_{j}-T_{\mathrm{pcm}})$. Comparing:
\begin{equation}
  \frac{K_{j}}{2\pi r_{i}\Ui}
  \;=\;
  \frac{1-e^{-\mathrm{NTU}_{j}}}{\mathrm{NTU}_{j}}
  \;\le\; 1 .
  \label{eq:K-ratio}
\end{equation}
The two limits say what the difference means:

\begin{description}[leftmargin=1.9em,itemsep=3pt]
\item[$\mathrm{NTU}\to0$:] the ratio $\to1$. The fluid barely cools across the
  segment, $T(z)\approx T_{j}$ throughout, and the naive closure is right.
\item[$\mathrm{NTU}\to\infty$:] $K\to \md c_{p}/\Delta z$, a finite limit set
  by the \emph{fluid}, not the wall. No matter how good the borehole is, a
  segment cannot absorb more than the stream can carry to it. The naive closure
  has no such limit and would let $T_{j+1}$ fall below $T_{\mathrm{pcm}}$ ---
  heat flowing from cold to hot, at a rate that grows without bound as $\Ui$
  improves.
\end{description}

\noindent
At this design point $\mathrm{NTU}$ runs from \num{0.083} to \num{0.602} over
the charge, so Eq.~\eqref{eq:K-ratio} lies between \num{0.96} and \num{0.75}:
\textbf{using $\Ui P_{i}$ would overstate the conductance by \SI{4}{\percent}
to \SI{33}{\percent}}, worst where the melt layer is thinnest and the borehole
best. The error is largest exactly where the model is most active, which is why
it cannot be waved away as second order.

\paragraph{Two properties inherited from this derivation.}

\begin{description}[leftmargin=1.9em,itemsep=3pt]
\item[Closure is structural.] The code does not evaluate
  Eq.~\eqref{eq:K-derived} and Eq.~\eqref{eq:pcm-balance} independently. It
  integrates Eq.~\eqref{eq:pcm-balance} over the step and \emph{reports}
  $q'_{j}=(E'^{\,n+1}-E'^{\,n})/\Delta t$, then drops the fluid temperature by
  exactly that amount, $T_{j+1}=T_{j}-q'_{j}\Delta z/(\md c_{p})$. Whatever the
  PCM gained, the fluid lost, to machine precision, with no possibility of
  drift.
\item[$K$ is frozen over a time step, $T_{\mathrm{pcm}}$ is not.] $K_{j}$ is
  evaluated once per step from the melt thickness at the start of it, whereas
  $T_{\mathrm{pcm}}$ is advanced exactly inside the step by
  Eq.~\eqref{eq:enth-exact}. The consequence is that $T_{j+1}$ equals the
  $\varepsilon$--NTU result of Eq.~\eqref{eq:ntu-derived} only while the segment
  is on the latent plateau, where $T_{\mathrm{pcm}}=\Tm$ is genuinely constant.
  In the sensible branches the two differ at second order in $\Delta t$. This
  is a deliberate choice: it keeps closure exact, which matters more than
  matching an $\varepsilon$--NTU relation that is itself only quasi-steady.
\end{description}

Two properties of Eq.~\eqref{eq:enth-ode} drive everything that follows:

\begin{itemize}[leftmargin=1.5em,itemsep=2pt]
\item It is \textbf{nonlinear}, because $T_{\mathrm{pcm}}(E')$ is piecewise.
\item It is \textbf{linear inside each branch}, because each piece of
  Eq.~\eqref{eq:enth-curve} is affine in $E'$.
\end{itemize}

\noindent
The second property is what makes an exact integration possible, and
\S\ref{sec:front:integration} takes it up. But first, why an exact integration
is not optional.

\subsubsection{Why explicit stepping fails --- and why nobody noticed until v0.4}
\label{sec:front:integration}

Look again at the three branches. On the plateau, $T_{\mathrm{pcm}}=\Tm$ is
\emph{constant}: it does not respond to $E'$ at all. The segment therefore
behaves as though it had infinite heat capacity, Eq.~\eqref{eq:enth-ode} reduces
to $\mathrm{d}E'/\mathrm{d}t = \text{constant}$, and an explicit Euler step is
not merely stable but \emph{exact}.

This is precisely the regime Formulations A and B live in permanently. Explicit
stepping was safe in v0.1 and v0.2 for a structural reason, not a lucky one ---
and that is why the time grid was never examined.

Add the sensible branches and the capacity becomes finite. Equation
\eqref{eq:enth-ode} is now a relaxation equation with time constant $C/K$, and
explicit Euler requires $\Delta t < C/K$. At the design point:

\begin{center}
\begin{tabular}{@{}l r r@{}}
\toprule
 & early ($t=\SI{1}{\second}$) & late ($t\to t_{ch}$) \\
\midrule
$\mathrm{NTU}$, segment 1        & \num{0.592}  & \num{0.083} \\
$K$ \si{[\watt\per\meter\per\kelvin]} & \num{64.3} & \num{11.5} \\
$C_{l}/K$ \si{[\second]}         & \num{197}    & \num{1106}  \\
time-grid step $\Delta t$ \si{[\second]} & \num{0.31} & \num{8491} \\
\midrule
$\Delta t\,K/C_{l}$              & \num{0.002}  & \textbf{\num{7.7}} \\
\bottomrule
\end{tabular}
\end{center}

\noindent
The grid is logarithmic, so it is fine at the start and violates the limit by
almost an order of magnitude at the end. The first implementation duly
diverged: the oscillation fed back through the fluid temperature and drove the
secondary fluid to \SI{261}{\kelvin}, some \SI{170}{\kelvin} below anything
physical.

This is a good failure to have seen once. It was loud, it was immediate, and it
was diagnosable. The dangerous version of the same error is a step ratio of
$1.5$ rather than $7.7$: still unstable in principle, but slowly enough to look
like a plausible answer.

\subsubsection{Branchwise-exact integration}

Rather than shrink $\Delta t$ --- which would multiply the cost of every run by
$\sim\!10^{3}$ --- exploit the linearity within each branch and integrate in
closed form. Holding the fluid temperature $T_{j}$ fixed across the step:

\begin{equation}
  \underbrace{E'(t+\Delta t) = E' + K\left(T_{j}-\Tm\right)\Delta t}
             _{\text{plateau: linear}},
  \qquad
  \underbrace{E'(t+\Delta t) = E'_{\mathrm{eq}}
    + \bigl(E'-E'_{\mathrm{eq}}\bigr)e^{-K\Delta t/C}}
             _{\text{sensible: exponential relaxation}},
  \label{eq:enth-exact}
\end{equation}

\noindent
with the equilibrium enthalpy and the relevant capacity
\begin{equation}
  \left(C,\;E'_{\mathrm{eq}}\right) =
  \begin{cases}
    \left(C_{s},\; C_{s}\left(T_{j}-\Tm\right)\right), & E' < 0,\\[3pt]
    \left(C_{l},\; E'_{\mathrm{lat}} + C_{l}\left(T_{j}-\Tm\right)\right),
      & E' > E'_{\mathrm{lat}}.
  \end{cases}
\end{equation}

$E'_{\mathrm{eq}}$ is the enthalpy at which the PCM would sit in equilibrium
with the fluid --- the state the segment relaxes toward and never quite reaches.
The exponential is unconditionally stable and monotone: no step size, however
large, can overshoot it. That is the property explicit Euler lacks.

\paragraph{Crossing a branch boundary.}
A step may start in one branch and end in another. The step is then split at the
crossing. Let $E'_{b}$ be the boundary being approached ($0$ or
$E'_{\mathrm{lat}}$). The crossing time is
\begin{equation}
  t^{*} =
  \begin{dcases}
    \frac{E'_{b}-E'}{K\left(T_{j}-\Tm\right)}, & \text{from the plateau},\\[8pt]
    \frac{C}{K}\,
      \ln\!\left(\frac{E'-E'_{\mathrm{eq}}}{E'_{b}-E'_{\mathrm{eq}}}\right),
      & \text{from a sensible branch},
  \end{dcases}
  \label{eq:crossing}
\end{equation}
and if $t^{*} < \Delta t$ the state is advanced to $E'_{b}$, the remaining
$\Delta t - t^{*}$ is integrated in the adjacent branch, and the test is
repeated. At most a handful of crossings can occur in one step, so the loop is
bounded.

The heat reported to the fluid is then
\begin{equation}
  q'_{j} = \frac{E'^{\,n+1}-E'^{\,n}}{\Delta t},
  \label{eq:qfromE}
\end{equation}
\emph{by definition}, not by a separate flux evaluation. Whatever the segment
gained, the fluid lost. Closure is exact to machine precision and cannot drift.

\begin{worked}[A bug worth studying: the sticky boundary]
The branch test in the code is written with closed intervals,
$0 \le E' \le E'_{\mathrm{lat}}$, so a state landing \emph{exactly} on
$E'_{\mathrm{lat}}$ is still classified as two-phase. Follow the logic: the
plateau branch computes $t^{*} = (E'_{\mathrm{lat}} - E')/\text{rate} = 0$,
finds $t^{*} < \Delta t$, advances the state to $E'_{\mathrm{lat}}$ ---
where it already is --- and subtracts zero from the remaining time. The next
iteration does the same. The segment is pinned at the boundary for ever and
never enters the liquid branch.

The symptom was not a crash. It was every segment reporting a melt fraction of
exactly $1.000$ and a superheat of exactly zero --- results that look
\emph{excellent}. The tell was a single diagnostic line printing ``sensible
energy stored: \num{0.000}'', a quantity that had no business being identically
zero.

The fix is to step just past the boundary rather than onto it,
\[
  E' \leftarrow E'_{\mathrm{lat}} + \epsilon,
  \qquad \epsilon = 10^{-9}\max\!\left(E'_{\mathrm{lat}},1\right),
\]
so the next iteration selects the liquid branch. The general lesson: when a
solver switches on the sign of a quantity, decide deliberately which side of the
boundary the equality belongs to, and make sure the state cannot come to rest
on it.
\end{worked}

\paragraph{Verification.}
The scheme was checked against a brute-force explicit integration of
Eq.~\eqref{eq:enth-ode} using \num{200000} substeps, over a charge and a
discharge, including steps that cross branch boundaries. Agreement is
$3\times10^{-13}$ relative --- machine precision. This is a genuine
verification: the two calculations share no code path beyond
Eq.~\eqref{eq:enth-curve}.

\subsubsection{Recovering the melt-layer thickness}

The resistance network needs $\delta$, not $\Amelt$. Inverting
Eq.~\eqref{eq:area-fins} --- which excludes the fin metal from the melted PCM
--- gives $\delta$ as described in \S\ref{sec:front:finmetal}, and the network
of \S\ref{sec:network} is otherwise unchanged in form.

One thing \emph{is} changed: the driving temperature. In Formulations~A and~B
the outer boundary of the resistance chain sits at $\Tm$. In Formulation~C it
sits at $T_{\mathrm{pcm}}$, which equals $\Tm$ only while some solid remains.
This single substitution is the mechanism of the next subsection.

\subsubsection{Sensible heat is self-levelling}
\label{sec:front:levelling}

Here is why a \SI{4.7}{\percent} energy term produces a first-order change in
the answer.

Consider again the well-head segment that melts first. Under Formulation~B it
finished melting and then continued to absorb heat at the full driving
difference $T_{j}-\Tm$, because its boundary temperature was pinned at $\Tm$
regardless of state. It monopolised the fluid's heat. The segments downstream,
which had not finished, received whatever was left.

Under Formulation~C the same segment starts to superheat the instant its solid
is gone. $T_{\mathrm{pcm}}$ rises, $T_{j}-T_{\mathrm{pcm}}$ collapses, its
demand falls away, and the heat continues down the well to segments that still
have solid to melt. The mechanism is a negative feedback, and it is automatic:
nothing in the code distributes the heat, and no parameter was tuned.

The effect on the distribution of melt is large:

\begin{center}
\begin{tabular}{@{}l l r@{}}
\toprule
formulation & $\varepsilon_{\mathrm{local}}$ at ten stations & spread \\
\midrule
B (latent only) & 1.78\; 0.73\; 0.69\; 0.71\; 0.77\; 0.85\; 0.94\; 1.03\; 1.12\; 1.23 & 1.321 \\
C (enthalpy)    & 1.00\; 1.00\; 1.00\; 1.00\; 1.00\; 1.00\; 1.00\; 1.00\; 1.00\; 1.00 & \textbf{0.078} \\
\bottomrule
\end{tabular}
\end{center}

\noindent
where the spread is
$\max_{j}\varepsilon_{\mathrm{local}}-\min_{j}\varepsilon_{\mathrm{local}}$ over
the \num{100} segments at end of charge. Two separate things improved. The
values above $1$ are gone, because Eq.~\eqref{eq:enth-curve} makes them
impossible. The values well below $1$ are gone too, and that is the physical
result: PCM that never melts is \emph{dead volume} --- it occupies borehole,
costs money, and carries no energy through the round trip. Formulation~C says
there is far less of it than Formulation~B implied.

\subsubsection{What Formulation C still does not do}

Three limitations, stated plainly because they bound how far the result can be
pushed.

\begin{description}[leftmargin=1.7em,itemsep=3pt]
\item[One temperature per segment.] $E'$ is a lumped quantity: all the PCM in a
  segment is at a single $T_{\mathrm{pcm}}$. There is no radial resolution
  \emph{within} a phase, so superheat appears instantaneously and uniformly
  across the whole liquid annulus rather than diffusing outward from the tube.
  The model therefore reaches its superheated state sooner than the real
  material would, and the self-levelling above is an upper bound on the effect.
\item[H1 is inherited, not repaired.] The quasi-steady melt layer of
  \S\ref{sec:front:legacy} is still assumed: the conduction profile through the
  liquid is taken as instantaneously steady. This is justified by
  $\mathrm{Ste}=\num{0.047}$ and is a good approximation here, but it is an
  approximation, and it is the same one in all three formulations.
\item[H3 is now \emph{saturated}.] Because $\varepsilon_{\mathrm{local}}\to1$
  over essentially the whole well, $\delta \to \delta_{\mathrm{merge}}$
  everywhere --- the melt fronts of neighbouring legs are touching. The annular
  conduction resistance is being evaluated at the exact limit of its validity.
  \S\ref{sec:defects:h3} takes this up, because it now carries a design
  decision.
\end{description}

Formulation~C is nonetheless the first of the three that is \emph{closed}: every
joule the resistance network delivers has a place to go, and the melt fraction
cannot leave $[0,1]$. Whatever is wrong with the answer is now wrong in the heat
transfer, not in the bookkeeping.

\emph{Source:} \code{pcm\_capacities}, \code{pcm\_state},
\code{advance\_segment}, \code{march\_h}.

%---------------------------------------------------------------------
\subsection{Latent inventory and the melted fraction}
\label{sec:front:inventory}

The melted volume per tube and the melted fraction of the borehole inventory are
\begin{equation}
  \Vmelt \;=\; \sum_{j=1}^{n_{s}} \Amelt^{(j)}\,\Delta z,
  \qquad
  \varepsilon_{\mathrm{PCM}}
  \;=\;
  \frac{\Vmelt\, n_{t}}{V_{\mathrm{well}}},
  \label{eq:eps}
\end{equation}
with $V_{\mathrm{well}}$ from Eq.~\eqref{eq:Vwell} and the multiplier $n_{t}$,
not $2n_{t}$, for the reason given after Eq.~\eqref{eq:Ltube}.

The density $\rhom$ in Eqs.~\eqref{eq:marched}--\eqref{eq:qeff} is the
\emph{same in both directions}, and is taken as the solid density $\rho_{s}$.
The thermal conductivity $\km$ by contrast \emph{is} directional, $k_{l}$ while
melting and $k_{s}$ while freezing, because the layer adjacent to the tube is
liquid in the first case and solid in the second.

The asymmetry is deliberate. Making the latent inventory directional as well
would bank energy at $\rho_{l}$ and return it at $\rho_{s}$, yielding
$\rho_{s}/\rho_{l}=1.069$ times more energy out than in --- a \SI{6.9}{\percent}
round-trip gain from bookkeeping alone. Carrying the front as a
\emph{solid-equivalent} area instead makes $\varepsilon_{\mathrm{PCM}}$ a true
melted \emph{mass} fraction and makes stored energy equal to available energy.
The price is that the \SI{6.9}{\percent} volume change on melting is not
represented: the liquid layer is about \SI{2.3}{\percent} thicker in radius than
Eq.~\eqref{eq:area-fins} reports, so Eq.~\eqref{eq:he} mildly under-predicts the
melt-layer resistance.

\emph{Source:} \code{thums\_core/stefan.py}; \code{config.PCM.rho\_latent}.
